\documentclass[11pt,a4paper]{article}
\usepackage[utf8]{inputenc}
\usepackage[vietnamese]{babel}
\usepackage{amsmath,amssymb}
\usepackage{geometry}
\geometry{a4paper, total={180mm,260mm}, left=15mm, right=15mm, top=18mm, bottom=18mm}

\usepackage{tikz}
\usetikzlibrary{shapes.geometric, arrows.meta, positioning, calc}
\tikzset{
  arr/.style={-{Stealth[scale=0.8]}},
  every node/.append style={align=center}
}
\usepackage{tcolorbox}
\tcbuselibrary{skins,breakable}
\usepackage{fontawesome}
\usepackage{fancyhdr}
\usepackage{enumitem}

\pagestyle{fancy}
\fancyhf{}
\renewcommand{\headrulewidth}{0.5pt}
\renewcommand{\footrulewidth}{0.5pt}
\fancyhead[L]{\small\bfseries\color{violet!85!black}\faBook\ SỔ TAY KIẾN THỨC TOÁN 2}
\fancyhead[R]{\small\bfseries\color{violet!85!black}Lớp Toán Thầy Lĩnh -- loptoanthaylinh.vn}
\fancyfoot[C]{\small\color{gray}-- Sổ tay kiến thức: Bài 39. Bảng nhân 2 --}

\setlength{\parindent}{0pt}
\setlength{\parskip}{0pt}

\begin{document}

% --- BANNER BÀI HỌC ---
\begin{tcolorbox}[
  enhanced,
  colback=violet!8!white,
  colframe=violet!70!black,
  arc=3mm, boxrule=1.2pt,
  left=3mm, right=3mm, top=2mm, bottom=2mm,
  halign=center
]
  {\LARGE\bfseries\color{violet!90!black}\faBookmark\ BÀI 39: BẢNG NHÂN 2}\\[1mm]
  {\itshape\footnotesize\color{black!75} Hình thành và ghi nhớ bảng cửu chương nhân 2; quy luật đếm thêm 2 trên 1 trang A4}
\end{tcolorbox}

\vspace{-2mm}

\noindent
\begin{minipage}[t]{0.49\textwidth}\vspace{0pt}
\begin{tcolorbox}[
  enhanced,
  colback=violet!4!white,
  colframe=violet!75!black,
  arc=2mm, boxrule=0.8pt,
  left=2mm, right=2mm, top=1.8mm, bottom=1.8mm,
  title={\footnotesize\bfseries\color{white}\faKey\ 1. Bảng nhân 2 chuẩn mực}
]
{\footnotesize
\begin{tabular}{l@{\quad}l}
  $2 \times 1 = \mathbf{2}$ & $2 \times 6 = \mathbf{12}$ \\
  $2 \times 2 = \mathbf{4}$ & $2 \times 7 = \mathbf{14}$ \\
  $2 \times 3 = \mathbf{6}$ & $2 \times 8 = \mathbf{16}$ \\
  $2 \times 4 = \mathbf{8}$ & $2 \times 9 = \mathbf{18}$ \\
  $2 \times 5 = \mathbf{10}$ & $2 \times 10 = \mathbf{20}$
\end{tabular}
\vspace{1mm}
\begin{itemize}\setlength{\itemsep}{0.5pt}\setlength{\leftskip}{-2.5mm}
  \item Thừa số thứ nhất đều là $2$.
  \item Thừa số thứ hai tăng dần từ $1$ đến $10$.
  \item \textbf{Tích} là dãy số chẵn tăng dần: $2, 4, 6, 8, \dots, 20$.
\end{itemize}
}
\end{tcolorbox}
\end{minipage}
\hfill
\begin{minipage}[t]{0.49\textwidth}\vspace{0pt}
\begin{tcolorbox}[
  enhanced,
  colback=teal!4!white,
  colframe=teal!75!black,
  arc=2mm, boxrule=0.8pt,
  left=2mm, right=2mm, top=1.8mm, bottom=1.8mm,
  title={\footnotesize\bfseries\color{white}\faSliders\ 2. Quy luật đếm thêm 2}
]
{\footnotesize
\begin{itemize}\setlength{\itemsep}{1pt}\setlength{\leftskip}{-2.5mm}
  \item Trong bảng nhân 2, tích liền sau luôn \textbf{hơn tích liền trước 2 đơn vị}.
  $$\text{Tích sau} = \text{Tích trước} + 2$$
  \item \textbf{Mẹo nhớ nhanh:}\\
  Nếu quên $2 \times 7$, ta nhớ $2 \times 6 = 12$ rồi lấy $12 + 2 = 14$.
  \item Dãy số cách đều 2: $2, 4, 6, 8, 10, 12, 14, 16, 18, 20$.
\end{itemize}
}
\end{tcolorbox}
\end{minipage}

\vspace{-1.5mm}
\begin{tcolorbox}[
  enhanced,
  colback=orange!4!white,
  colframe=orange!80!black,
  arc=2mm, boxrule=0.8pt,
  left=2mm, right=2mm, top=1.5mm, bottom=1.5mm,
  title={\footnotesize\bfseries\color{white}\faSitemap\ 3. Sơ đồ tư duy: Bậc thang dãy tích Bảng nhân 2}
]
\centering
\begin{tikzpicture}[x=0.9cm, y=0.6cm]
  \foreach \i/\val in {1/2, 2/4, 3/6, 4/8, 5/10, 6/12, 7/14, 8/16, 9/18, 10/20} {
    \node[circle, draw=violet!80!black, fill=violet!15, inner sep=2pt, font=\tiny\bfseries] at (\i*1.1, 0) {\val};
    \ifnum\i<10
      \draw[->, thick, orange!80!black] (\i*1.1+0.3, 0.2) to[bend left=40] node[above, font=\tiny] {+2} (\i*1.1+0.8, 0.2);
    \fi
  }
\end{tikzpicture}
\end{tcolorbox}

\vspace{-1.5mm}
\noindent
\begin{minipage}[t]{0.49\textwidth}\vspace{0pt}
\begin{tcolorbox}[
  enhanced,
  colback=green!4!white,
  colframe=green!60!black,
  arc=2mm, boxrule=0.8pt,
  left=2mm, right=2mm, top=1.8mm, bottom=1.8mm,
  title={\footnotesize\bfseries\color{white}\faCheckCircle\ 4. Ví dụ thực hành chuẩn mực}
]
{\footnotesize
\textbf{1. Tính nhẩm:}\\
$2 \times 4 = \mathbf{8}$; $2 \times 8 = \mathbf{16}$; $2 \times 9 = \mathbf{18}$.
\textbf{2. Bài toán lời văn:}\\
Mỗi đôi đũa có $2$ chiếc đũa. Hỏi $7$ đôi đũa có tất cả bao nhiêu chiếc đũa?\\
\textit{Bài giải:}\\
$7$ đôi đũa có số chiếc đũa là:
$$2 \times 7 = 14\text{ (chiếc)}$$
Đáp số: $14$ chiếc đũa.
}
\end{tcolorbox}
\end{minipage}
\hfill
\begin{minipage}[t]{0.49\textwidth}\vspace{0pt}
\begin{tcolorbox}[
  enhanced,
  colback=red!4!white,
  colframe=red!75!black,
  arc=2mm, boxrule=0.8pt,
  left=2mm, right=2mm, top=1.8mm, bottom=1.8mm,
  title={\footnotesize\bfseries\color{white}\faExclamationTriangle\ 5. Mẹo tránh bẫy trắc nghiệm}
]
{\footnotesize
\begin{itemize}\setlength{\itemsep}{1pt}\setlength{\leftskip}{-2.5mm}
  \item \textbf{Phân biệt $2 \times 3$ và $2 + 3$:}\\
  $2 \times 3 = 6$, còn $2 + 3 = 5$. Tránh nhìn nhầm dấu nhân thành dấu cộng!
  \item Tất cả các tích trong bảng nhân 2 đều là \textbf{số chẵn} tận cùng là $0, 2, 4, 6, 8$.
\end{itemize}
}
\end{tcolorbox}
\end{minipage}


\end{document}
